% This is samplepaper.tex, a sample chapter demonstrating the
% LLNCS macro package for Springer Computer Science proceedings;
% Version 2.20 of 2017/10/04
%
\documentclass[runningheads]{llncs}
%
\usepackage{lipsum}  
\usepackage{graphicx}
% Used for displaying a sample figure. If possible, figure files should
% be included in EPS format.
%
% If you use the hyperref package, please uncomment the following line
% to display URLs in blue roman font according to Springer's eBook style:
% \renewcommand\UrlFont{\color{blue}\rmfamily}

\begin{document}
%
\title{SASYR Abstract Title}
%
%\titlerunning{Abbreviated paper title}
% If the paper title is too long for the running head, you can set
% an abbreviated paper title here
%
\author{First Author\inst{1}\orcidID{0000-1111-2222-3333} \and
Second Author\inst{2,3}\orcidID{1111-2222-3333-4444} \and
Third Author\inst{3}\orcidID{2222--3333-4444-5555}}
%
\authorrunning{F. Author et al.}
% First names are abbreviated in the running head.
% If there are more than two authors, 'et al.' is used.
%
\institute{Instituto Politécnico de Viana do Castelo, IPVC, Portugal\\
\email{author@ipvc.pt}\\
\and
CeDRI, Instituto Politécnico de Bragança, Portugal\\
\email{author@ipb.pt}
\and
Instituto Politécnico do Cávado e do Ave, IPCA, Portugal\\
\email{author@ipca.pt}}
%
\maketitle              % typeset the header of the contribution
%
%
%
\section*{Abstract}
This is the abstract template to be used for SASYR. Please keep your abstract within the length given. The abstract should be in English. You cannot include images or tables. Just include text or mathematical symbols. The References are not mandatory.

\lipsum[1-1]


oindent Displayed equations are centered and set on a separate
line.
\begin{equation}
x + y = z
\end{equation}



oindent Displayed equations are centered and set on a separate
line.
\begin{equation}
x + y = z
\end{equation}




oindent Displayed equations are centered and set on a separate
line.
\begin{equation}
x + y = z
\end{equation}


% the environments 'definition', 'lemma', 'proposition', 'corollary',
% 'remark', and 'example' are defined in the LLNCS documentclass as well.
%

For citations of references,
In \cite{einstein} Einstein....


\keywords{First keyword  \and Second keyword \and Another keyword.}
%
% ---- Bibliography ----
%
% BibTeX users should specify bibliography style 'splncs04'.
% References will then be sorted and formatted in the correct style.
%
\bibliographystyle{refs-style}
\bibliography{refs}
%

\end{document}
